\section{Related Work} \label{sec:related}

\paragraph{Data Reuse in LLM Pre-Training.} Empirically, there is a long debate over the effect of data reuse in LLM pre-training. Some works~\citep{lee2021deduplicating,hoffmann2022training,hernandez2022scaling,wang2023data} suggested it may be harmful, while some works~\citep{taylor2022galactica} reported the benefit of data reuse when the number of epochs is small ($K \le 4$).
% \kaifengclaw{Typo/consistency: 'some work' -> 'some works'; and 'data reusing' is better unified as 'data reuse'.} 
\citet{xue2023repeat} then discovered a degradation phenomenon in multi-epoch training and investigated relevant factors and regularization methods to tackle it. \citet{muennighoff2023scaling} trained LLMs under different configurations and also found that reusing data is as good as using fresh data in the first few epochs. Yet, as the number of epochs increases, the returns for repetitions diminish. In our work, from a theoretical perspective, we rigorously analyze the effect of data reuse using non-asymptotic techniques, and we define and calculate the effective reuse rate under two cases, shedding light on the theoretical understanding of data reusing in LLM pre-training.
% \kaifengclaw{Tense/style: keep consistent tense in related work (e.g., present tense: 'analyze', 'define', 'calculate').}

\paragraph{Comparison with \citet{lin2025improved}.} A recent study on linear regression with data reusing~\citep{lin2025improved} is among the most relevant to our results. They showed that when the number of epochs is relatively small (smaller than some power of the dataset size), the order of loss remains the same as one pass SGD for the same iterations, which aligns with our results. However, their results only imply that $E(K,N) = \Theta(K)$ for small $K$, while our analysis directly gives the explicit loss characterization with $o(1)$ relative error bound and a more exact description of the effective reuse rate, which reflects the data reusing scaling behaviour. Our analysis is across various problem setups, and further shows the general scaling trend of data reusing under different problem setups.

%\tingkai{Read all the papers and separate the papers; emphasize that some paper consider the linear regression setting; more papers to cite}
\iffalse
\paragraph{Data Reuse in LLM Pre-Training.} 
% As fresh data is becoming scarce \citet{villalobos2024run}, data reuse is getting more and more attention and remains a research hotspot empirically and theoretically. Some works have focused on the downside of data reuse, for example, 

%There are many empirical discussions about the effect of data reusing in Large Language Model (LLM) training. 
Some previous works suggested that data reuse for LLM training may be harmful~\citep{lee2021deduplicating,hoffmann2022training,hernandez2022scaling,wang2023data}. However, \citet{taylor2022galactica} observed that reusing pre-training data up to $4$ epochs for a $120$B model has a positive impact. \citet{xue2023repeat} then discovered a degradation phenomenon in multi-epoch training and investigated relevant factors and regularization methods to tackle it. \citet{muennighoff2023scaling} trained large language models under different configurations and found that reusing previous data is as good as using new fresh data in the first few epochs. Yet, as the number of epochs increases, the returns for repetitions diminish. 
% \haodong{check empirical papers: some LLM empirical papers show that data repeating and match fresh data only when the number of epochs is very small. Some theory - > better for repeating, but observed on some synthetic setting}
Besides the real-world LLM pre-training regime, many works also reported the improvement of data reusing under synthetic settings empirically~\citep{charton2024emergent,kazdan2024collapse} or theoretically~\citep{zucchet2025emergence,dandi2024benefitsreusingbatchesgradient,arnaboldi2025repetitaiuvantdatarepetition}.  
% \citet{charton2024emergent}studied the performances of transformers on three types of math problems and concluded that data repetition is necessary. 
% Specifically, on the theoretical side, a line of work developed tools in mean field theory and gave results on the limitations of data reuse under some synthetic settings. \citet{dandi2024benefitsreusingbatchesgradient} considered a function-learning problem with multi-pass gradient descent that reuses the batches multiple times. This optimizer performed efficiently and overcame the constraints given by the information exponent \citep{BenArous2021} and leap exponent \citep{abbe2023sgd}, even for functions that does not satisfy the staircase property \citep{abbe2024mergedstaircasepropertynecessarynearly}. 
% Furthermore, \citet{arnaboldi2025repetitaiuvantdatarepetition} investigated a modification of single-pass gradient descent where data could be reused twice. They observed a drastic improvement in computational efficiency. 
In our work, we rigorously analyzed the effect of data reuse using non-asymptotic techniques, and we defined and calculated the effective reuse rate under two cases, shedding light on the theoretical understanding of data reusing in LLM pre-training.

% \paragraph{Empirical Scaling Laws.}
% Scaling laws that reveal the relationships between training loss and other factors such as model size, data size and compute budget have always been an extensive area of research. \citet{kaplan2020scaling} is the first to propose a power-law decay formula w.r.t. data size and model size in the context of training transformer-based large language models. Since then, literature have been focusing on fitting scaling laws under different settings ... . \citet{hoffmann2022training} put forward a Chinchilla law that is different from Kaplan scaling laws, and proposed methods to allocate data and model size under a given compute budget. \citet{muennighoff2023scaling} fitted a more-complex scaling law which takes the number of training epochs into consideration.
%\haodong{ rephrase the following scaling law reference}

\paragraph{Empirical Findings on Scaling Laws.} Scaling laws reveal the relationships between large-scale model training loss and various factors such as model size, data size, and compute budget. These laws were initially observed by \citet{hestness2017deep}, but gained significant influence through the work of \citet{kaplan2020scaling}, and have since been further developed in a series of studies \citep{henighan2020scaling, hoffmann2022training, zhai2022scaling, kadra2023power, aghajanyan2023scaling, muennighoff2023scaling, bi2024deepseek, shuai2024scaling, kumar2024scaling, tissue2024scaling, luo2025multi}. %Notably, \citet{hoffmann2022training} introduced the Chinchilla Law, which proposes an optimal scaling strategy under a fixed compute budget. Additionally, \citet{tissue2024scaling} and \citet{luo2025multi} explored scaling laws with learning rate schedules (LRSs) based on empirical observations. 
Notably, \citet{muennighoff2023scaling} further refined these models by incorporating the number of training epochs into a more complex scaling law, which empirically describes the effect of data reuse. In our work, we provide a theoretical analysis of how the effective reuse rate $E(K,N)$ relies on the epoch number $K$ and fresh data size $N$, highlighting the role of $N$ in the scaling behavior of $E(K,N)$, a factor that was overlooked in \citet{muennighoff2023scaling}. 

% \paragraph{Theory of Scaling laws.}
% There are many existing works that interpreted the mechanism behind the scaling laws ... . \citet{sharma2020neuralscalinglawdimension} considered a toy model and showed that $\cL \propto \frac{1}{N^{\frac{4}{d}}}$, where $d$ is the intrinsic dimension. \citet{Bahri_2024} and \citet{bordelon2024dynamicalmodelneuralscaling} both analyzed a linear teacher-student model setting with a power-decay parameter prior. They proposed theories for scaling laws based on statistical heruistics. \citet{bordelon2025featurelearningimproveneural} developed a solvable model of neural scaling laws and defined regimes according to task difficulties. Meanwhile, other works considered more complicated settings and gave more rigorous mathematical proofs. \citet{jain2024scalinglawslearningreal} studied how data from other sources affect the test error of models and their scaling laws. \citet{paquette202543phasescomputeoptimalneural} identified different phases for a PLRF model in the compute-limited, infinite-data scaling law regime, and they discovered different scaling exponents in different phases. \citet{lin2024scaling} proved a scaling law using a PLRF model with an online SGD scheme with a geometrically decaying stepsize non-asymptotically. However, these scaling law papers failed to consider the impact of data reuse. In contrast, we proved matching upper and lower risk bounds in two settings for multi-epoch training.
% \paragraph{Theory of Multi-Pass SGD.}
% Shuffling is a common practice in machine learning, however, theoretical literature on SGD with shuffling is limited. 

%\haodong{emphasize that the linear regression setting is attracting more and more interest}
\paragraph{Theoretical Explanations for Scaling Laws.} A series of studies \citep{sharma2020neuralscalinglawdimension,hutter2021learning,maloney2022solvable,wei2022more,jain2024scalinglawslearningreal,michaud2024quantization,nam2024exactly,atanasov2024scaling,dohmatob2024tale,bahri2024explaining,bordelon2024dynamical,lin2024scaling,paquette202543phasescomputeoptimalneural,bordelon2024feature,zhang2024does,ferbach2025dimension,li2025fsl} have sought to theoretically explain scaling laws from various perspectives. Among these, recent works \citep{bordelon2024dynamical,paquette202543phasescomputeoptimalneural,lin2024scaling,bordelon2024feature} have analyzed scaling laws by tracking the training dynamics of SGD through linear regression setup. Specifically, \citet{bordelon2024dynamical} investigated a full-batch gradient flow setup, while \citet{paquette202543phasescomputeoptimalneural} and \citet{bordelon2024feature} focused on online SGD with a sufficiently small constant learning rate. Additionally, \citet{lin2024scaling} studied a geometric decaying learning rate schedule (LRS) \citep{ge2019step,wu2022iterateriskboundssgd}. Recently, \citet{li2025fsl} proposed a functional scaling law that characterizes the loss dynamics for general LRSs. However, these scaling law studies did not account for the impact of data reuse. In contrast, our work examines the scaling behavior of multi-epoch SGD training within the context of a linear regression setup. 

\paragraph{Comparison with \citet{lin2025improved}.} A recent study on linear regression with data reusing~\citep{lin2025improved} is among the most relevant to our results. They showed that when the number of epochs is relatively small (some power of the dataset size), the order of loss remains the same as one pass SGD for the same iterations, which aligns with our results. However, their results only imply that $E(K,N) = \Theta(K)$ for small $K$, while our analysis directly gives the explicit loss characterization with $o(1)$ relative error bound and a more exact description of the effective reuse rate, which reflects the data reusing scaling behaviour. Our analysis is across various problem setups, and further shows the general scaling trend of data reusing under different problem setups.

% A recent related paper \citet{lin2025improved} studies the scaling law in the data reuse regime and showed that the loss under data reusing can improve the exponents in the scaling law for a class of regression problems. Their analysis considers the shuffling with replacement, and is hard to derive an explicit characterization of the data-resuse efficiency in general linear regression problems.\haodong{point out the defference of Theta and 1 + o(1);why they can iimporve but we cannot}

\paragraph{SGD Analysis in Linear Regression.} %\haodong{mention some scaling law papers} 
The analysis of SGD in linear regression has been extensively studied over the years, encompassing both one-pass and multi-epoch SGD. In the context of one-pass SGD, \citet{zou2021benign,meterez2025simplified} considered an SGD procedure with a constant step size and averaged iterates, offering a sharp risk bound in terms of the eigenvalues of the covariance matrix. \citet{pmlr-v139-gurbuzbalaban21a} examined one-pass SGD with batch size and proved that the distribution of
the SGD iterates will converge to a heavy-tailed
stationary distribution. \citet{zou2022benefitsimplicitregularizationsgd} compared the performance of SGD in the absence of ridge regression. \citet{wu2022iterateriskboundssgd} and \citet{wu2022powerlimitationpretrainingfinetuninglinear} studied SGD in linear regression under covariate shift. \citet{xia2024data} considered SGD updates with noisy gradient and analyzed the perfect deleted point problem. \citet{li2025understanding} considered SGD with exponential moving average in the linear regression setting. For multi-epoch SGD, \citet{lin2019optimalratesmultipassstochastic} examined a scenario in which gradients are sampled uniformly at random and mini-batches are allowed. They analyzed the effects of mini-batch size, number of epochs, and learning rate, carefully combining these parameters to achieve the optimal convergence rate. \citet{pillaudvivien2018statisticaloptimalitystochasticgradient} showed that while single-pass averaged SGD is optimal for a certain class of "easy" problems, multiple passes are required to achieve optimal prediction performance on a different class of "hard" problems, provided that an appropriate step size is chosen. In contrast to the matching upper and lower bounds derived by our theory, however, all the above works were only able to derive an upper bound for the loss.

% In one-pass SGD for linear regression, constant-step-size schedules with iterate averaging yield sharp risk bounds based on covariance eigenvalues (Zou et al., 2021; 2022) and have been studied under covariate shift (Wu et al., 2022). For multi-epoch SGD, optimal rates depend on mini‑batch size, epoch count, and learning rate (Lin et al., 2019), and while averaged single‑pass SGD is optimal for “easy” problems, “hard” problems require multiple passes with tuned step sizes, though only upper loss bounds were previously shown (Pillaud‑Vivien et al., 2018).

% \paragraph{SGD Analysis in Linear Regression.}
% One-pass SGD with constant step size and iterate averaging yields sharp eigenvalue-based risk bounds \citep{zou2021benign}, benefits without ridge regularization \citep{zou2022benefitsimplicitregularizationsgd}, and robustness under covariate shift \citep{wu2022iterateriskboundssgd,wu2022powerlimitationpretrainingfinetuninglinear}. Multi-epoch SGD achieves optimal convergence when mini‑batch size, epoch count, and learning rate are jointly tuned \citep{lin2019optimalratesmultipassstochastic}, and while single-pass averaging suffices for “easy” problems, “hard” problems require multiple passes—though only upper loss bounds are known \citep{pillaudvivien2018statisticaloptimalitystochasticgradient}.
\fi
